\section{Распознавание на устройстве: данные, модель и ESP-DL}\label{les:edge-ai}
\lessonintro
  {различать обучение и инференс и измерять качество}
  {уроки 6, 20, 26 и 35}
  {классификатор двух положений ручки на собственных данных}
  {потенциометр GPIO1; для расширения --- IMU или камера}

\subsection{Что даёт именно ESP32-S3}
LX7 имеет векторные инструкции для ускорения некоторых вычислений; PSRAM полезна для моделей
и изображений. Но произвольный цикл C++ не ускорится от установки библиотеки: нужны
поддерживаемые операции и оптимизированные ядра. ESP-DL --- библиотека Espressif для инференса.
Проверяйте совместимые модели, конвертер и версию ESP-IDF по конкретному выпуску.

Обучение настраивает параметры по примерам, инференс применяет их к новому входу.
Начнём с маленькой модели: среднее измерение для каждого из двух классов.
Она не является нейросетью и не использует векторное ускорение, но показывает весь путь
от данных к предсказанию, включая ошибки и неизвестный класс.

\subsection{Обучение на своих измерениях}
Определите классы 0 и 1 как два диапазона положения ручки. Подключите потенциометр как в уроке 6.
Поставьте ручку в первый диапазон и отправьте 0, затем во второй и отправьте 1.
На класс собирается 100 отсчётов; слегка меняйте положение внутри выбранного диапазона.

\codecap{Обучение двух центров и отказ от далёкого входа}
\codeinput{l37_classifier}

Нормализация приводит отсчёт к масштабу 0--1. Модель выбирает ближайший центр.
Если оба далеко, возвращает −1 («неизвестно»). Порог 0,12 учебный, подбирайте его
на отдельных данных. Расстояние --- не вероятность. После сброса центры исчезают:
сохранение обученной модели в NVS --- отдельная задача с проверкой версии формата.

\subsection{Как проверить честно}
После обучения заново установите ручку в тестовые положения: те же 100 отсчётов не являются
независимым тестом. Для жестов разделяйте целые записи и сеансы, а не соседние отсчёты движения:
иначе оценка будет завышена из-за похожих данных. Включите промежуточные положения и фон.
Составьте таблицу истинного и предсказанного класса, включая «неизвестно».
Считайте не только долю успехов, но и ложные срабатывания. Сравните с простым порогом:
модель оправдана, если улучшает решение конкретной задачи.

\subsection{Переход к IMU, камере и ESP-DL}
Для IMU выделяют окно, затем признаки: среднее, разброс, пики, энергию.
Для камеры фиксируют размер входа, порядок каналов и нормализацию, одинаковые при обучении
и применении. Квантование int8 может уменьшить память, но качество проверяют после преобразования.
Сквозной путь с настоящей нейросетью, квантованием и измерениями разобран в уроке 42. Выберите официальный пример ESP-DL для ESP32-S3, зафиксируйте версию и рекомендованный конвертер.
Не смешивайте произвольную модель с несовместимым форматом исполнения.

Память нужна не только весам: входы, промежуточные тензоры, кадры, стек и сеть тоже занимают место.
Измеряйте отдельно сбор данных, предобработку, инференс и реакцию.
Команда мотору требует подтверждения или нескольких устойчивых совпадений: единичная ошибка
распознавания не должна немедленно запускать движение.

\begin{trouble}
Один класс на всех входах: центры обучены рядом. Качество только дома: нет других условий в данных.
Тормозит камера: причина может быть в копировании и предобработке, а не модели.
\end{trouble}
\begin{summary}
Качество измеряют на новых данных вместе с ошибками, задержкой и памятью.
Аппаратная возможность ускорения не заменяет измерения конкретного приложения.
\end{summary}

\begin{selfcheck}
\begin{enumerate}
\item Чем обучение отличается от инференса?
\item Почему соседние отсчёты жеста нельзя случайно делить между обучением и тестом?
\item Почему размер модели меньше памяти приложения?
\end{enumerate}
\end{selfcheck}
\needspace{20\baselineskip}
\begin{tasks}
\begin{enumerate}
\item \lv{1} Опишите два класса так, чтобы другой ученик повторил опыт.
\item \lv{1} Обучите центры и проверьте по пять новых положений.
\item \lv{2} Проверьте десять промежуточных положений и настройте отказ.
\item \lv{2} Сохраните модель в NVS с номером формата.
\item \lv{3} Соберите жесты IMU из разных сеансов и сравните модель с порогами.
\item \lv{4} Запустите официальный пример ESP-DL и опишите версии, вход, память, задержку и матрицу ошибок.
\end{enumerate}
\end{tasks}

\paragraph{Документация для продолжения}\href{https://github.com/espressif/esp-dl}{Официальное руководство по теме}.
